\section*{Capítulo \ref{ch:b2:structure-determination} --- Determinação estrutural: RMN de \texorpdfstring{$^{13}$C}{13C}, EM e IV combinados}

\begin{solution}{exo:b2:structure-determination:1}
1,2-dimetilbenzeno: 4 (\ce{CH3}; C1/C2; C3/C6; C4/C5). 1,3-: 5 (\ce{CH3}; C1/C3; C2; C4/C6; C5).
1,4-: 3 (\ce{CH3}; C1/C4; C2/C3/C5/C6).
\end{solution}

\begin{solution}{exo:b2:structure-determination:2}
\ce{CH3CH(OH)CH2CH3}. DEPT-135: C1 (\ce{CH3}) para cima, C2 (CH) para cima, C3 (\ce{CH2}) para baixo,
C4 (\ce{CH3}) para cima. DEPT-90: só C2.
\end{solution}

\begin{solution}{exo:b2:structure-determination:3}
209: carbonila de cetona; 170: carbonila de éster (ou de ácido); 129: carbono \omterm{def:b2:huckel:aromatic}{aromático}; 64: carbono
ligado a um oxigênio; 14: \ce{CH3} alquila.
\end{solution}

\begin{solution}{exo:b2:structure-determination:4}
$(2 \times 8 + 2 - 8)/2 = 5$: um anel benzênico (4) e uma \ce{C=O} (1). Por exemplo o benzoato de metila,
\ce{C6H5COOCH3}, ou o ácido 4-metilbenzoico.
\end{solution}

\begin{solution}{exo:b2:structure-determination:5}
Uma insaturação, uma cetona (\qty{1715}{cm^{-1}}, raia em 209, quaternária); um \ce{CH2} e dois
\ce{CH3}, quatro carbonos diferentes: butan-2-ona, \ce{CH3COCH2CH3}.
\end{solution}

\begin{solution}{exo:b2:structure-determination:6}
Carbono: a pentan-2-ona mostra cinco raias, a pentan-3-ona, simétrica, três. Prótons: a pentan-2-ona tem
um singleto de três prótons (\ce{CH3CO}); a pentan-3-ona só um quarteto e um tripleto. \omterm{def:b2:mass-spec-atomic:spectrum}{Espectro de massas}:
a pentan-2-ona tem seu \omterm{def:b2:mass-spec-atomic:spectrum}{pico base} em 43 (\ce{CH3CO+}) e um íon de McLafferty em 58; a pentan-3-ona tem seu
\omterm{def:b2:mass-spec-atomic:spectrum}{pico base} em 57 (\ce{C2H5CO+}).
\end{solution}

\begin{solution}{exo:b2:structure-determination:7}
\ce{C4H8O2}, 88: o par quarteto--tripleto é um \ce{OCH2CH3} (o \ce{CH2} em 4.12 está no oxigênio), o
singleto um \ce{CH3CO}. Etanoato de etila, \ce{CH3COOCH2CH3}.
\end{solution}

\begin{solution}{exo:b2:structure-determination:8}
1,2-diclorobenzeno: três raias (C1/C2, C3/C6, C4/C5). 1,4-diclorobenzeno: duas (C1/C4, os quatro
CH). O DEPT-90 mostra duas raias CH para o primeiro, uma para o segundo.
\end{solution}

\begin{solution}{exo:b2:structure-determination:9}
C2 é um estereocentro. Substituir um ou outro próton de C3 por um deutério dá dois diastereoisômeros,
não enantiômeros: os dois prótons são diastereotópicos, em ambientes diferentes. Eles podem ter
deslocamentos diferentes e se acoplar entre si assim como com seus vizinhos, de modo que o
\ce{CH2} dá um multipleto complexo em vez de um simples quinteto.
\end{solution}

\begin{solution}{exo:b2:structure-determination:10}
Cinco insaturações; um anel monossubstituído (três raias CH e uma raia C em 137.0, cinco H \omterm{def:b2:huckel:aromatic}{aromáticos});
uma cetona (200.8, C); um grupo etila (\ce{CH2} 31.8, quarteto em 2.98 vizinho da carbonila, \ce{CH3}
8.3). 1-fenilpropan-1-ona, \ce{C6H5COCH2CH3}. A conjugação com o anel deslocaliza os elétrons $\pi$ da
\ce{C=O} e enfraquece a ligação: a banda cai abaixo dos \qty{1715}{cm^{-1}} de uma cetona saturada.
% ledger: ir:CO-ketone
\end{solution}

\begin{solution}{exo:b2:structure-determination:11}
Ácido pentanoico: a banda \ce{O-H} muito larga de 2500 a \qty{3300}{cm^{-1}} e um próton perto de 11--12
ppm. Butanoato de metila: um singleto de três prótons na faixa dos prótons de um carbono ligado a um
oxigênio (3.3--4.5). Propanoato de etila: dois quartetos, um no oxigênio e outro vizinho da carbonila.
Etanoato de propila: um singleto de três prótons vizinho da carbonila (2.0--2.4) e um tripleto no oxigênio.
% ledger: ir:OH-acid, nmrr:acid, nmrr:alcohol, nmrr:ketone
\end{solution}

\begin{solution}{exo:b2:structure-determination:12}
As duas raias desaparecem em \omterm{def:b2:structure-determination:dept}{DEPT}, portanto ambas são quaternárias, e o \omterm{def:b2:structure-determination:dept}{DEPT} sozinho não pode
distingui-las. Os deslocamentos podem: as carbonilas de éster ficam perto de 167--171 e os carbonos
\omterm{def:b2:huckel:aromatic}{aromáticos} entre 128 e 137 no gráfico deste capítulo. Um carbono de anel em 166.5 e uma carbonila em
130.6 ficariam ambos muito fora de suas classes; a atribuição correta é a inversa.
\end{solution}

\begin{solution}{pb:b2:structure-determination:1}
\textbf{1.} $3.1/31.7 \approx 9.8~\%$; $9.8/1.11 \approx 9$ carbonos.
\textbf{2.} Massa par: nenhum nitrogênio (ou dois).
\textbf{3.} $150 - 108 = 42$: \ce{H10O2} convém (10 + 32): \ce{C9H10O2}.
\textbf{4.} Dez carbonos dariam $M+1 \approx 11~\%$ de $M$, não 9.8~\%; e o IV mostra uma
carbonila, que exige um segundo oxigênio além da banda \ce{C-O}.
\textbf{5.} $(2 \times 9 + 2 - 10)/2 = 5$.
\textbf{6.} $108 + 10 \times 1.007825 + 2 \times 15.994915 \approx 150.0681$.
\textbf{7.} Um estiramento \ce{C=O} na posição de um éster saturado (perto de \qty{1735}{cm^{-1}}).
\textbf{8.} O estiramento \ce{C-O} do éster.
\textbf{9.} Ligações \ce{C-H} \omterm{def:b2:huckel:aromatic}{aromáticas} (ou de alceno).
\textbf{10.} Um grupo \ce{O-H}: nem álcool, nem ácido carboxílico.
\textbf{11.} Não: uma carbonila conjugada com um anel absorve em número de onda mais baixo; aqui ela fica
onde ficam os ésteres saturados, de modo que um grupo não conjugante a separa do anel.
\textbf{12.} 170.70: \ce{C=O} de éster; 136.14: carbono do anel que leva o substituinte; 128.56 e 128.24:
CH \omterm{def:b2:huckel:aromatic}{aromáticos}; 66.24: \ce{CH2} ligado ao oxigênio; 20.82: \ce{CH3}.
\textbf{13.} 66.24, o \ce{CH2}.
\textbf{14.} 170.70 e 136.14, os dois carbonos sem hidrogênio.
\textbf{15.} 128.56 e 128.24.
\textbf{16.} Quatro: o carbono que leva o substituinte, os dois orto, os dois meta e o CH para. Esperam-se
três tipos de CH; dois deles têm deslocamentos próximos demais para serem resolvidos neste campo, e
compartilham uma raia.
\textbf{17.} Não necessariamente: alturas não são contagens (\cref{prop:b2:structure-determination:intensities}),
embora aqui a raia possa bem conter dois tipos de CH sobrepostos.
\textbf{18.} 7.33: os cinco prótons \omterm{def:b2:huckel:aromatic}{aromáticos}; 5.09: o \ce{OCH2}; 2.06: o \ce{CH3CO}.
\textbf{19.} Seus átomos vizinhos (o carbono do anel e o oxigênio para o \ce{CH2}, o carbono carbonílico
para o \ce{CH3}) não levam hidrogênio.
\textbf{20.} Ele está ligado ao oxigênio eletronegativo do éster e ao anel: ambos o desblindam.
\textbf{21.} \ce{C6H5CH2OCOCH3}, etanoato de benzila.
\textbf{22.} No isômero, o \ce{CH2} não está ligado a um oxigênio, de modo que não poderia aparecer em
5.09, e a metila, no oxigênio, apareceria na faixa 3.3--4.5, não em 2.06 vizinha de uma carbonila.
% ledger: nmrr:alcohol, nmrr:ketone
\textbf{23.} 108: perda de 42 (ceteno, \ce{CH2=C=O}) por rearranjo, o cátion radical de
\ce{C7H8O}; 91: \ce{C7H7+}, o cátion benzila; 43: \ce{CH3CO+}.
\textbf{24.} Sete tipos de carbono (\ce{C=O}, o carbono substituído do anel, CH orto, meta e para,
\ce{CH2}, \ce{CH3}): \textbf{sete raias, uma delas (a do \ce{CH2}) apontando para baixo} em DEPT-135.
\end{solution}
